\documentclass[9pt,a4paper]{extarticle}

% ======================
% Codificação, idioma e layout
% ======================
\usepackage[T1]{fontenc}
\usepackage[utf8]{inputenc}
\usepackage[brazil]{babel}
\usepackage{geometry}
\geometry{margin=2cm}
\usepackage{parskip}
\usepackage{setspace}
\onehalfspacing
\usepackage{microtype}

% ======================
% Matemática e símbolos
% ======================
\usepackage{amsmath, amssymb, amsthm}
\usepackage{bm}

% ======================
% Figuras, tabelas e cores (se quiser adicionar desenhos/contornos)
% ======================
\usepackage{graphicx}
\usepackage{float}
\usepackage{caption}
\usepackage{subcaption}
\usepackage{xcolor}
\usepackage{booktabs}
\usepackage{array}
\usepackage{tabularx}

% ======================
% Listas e hyperlinks
% ======================
\usepackage{enumitem}
\setlist{nosep}
\usepackage{hyperref}
\hypersetup{
  colorlinks=true,
  linkcolor=black,
  urlcolor=blue
}

% ======================
% Título
% ======================
\title{\vspace{-1cm}Roteiro de Aula\\Otimização Não Linear\\Métodos de Busca Indireta (Descent Methods)}
\author{Prof. Rodrigo Silva}
\date{}

\begin{document}

\maketitle

\section*{Objetivos da aula}

Ao final desta aula, o estudante deverá ser capaz de:
\begin{itemize}
    \item Definir o \textbf{gradiente} de uma função $f:\mathbb{R}^n \to \mathbb{R}$ e interpretar seu significado geométrico.
    \item Explicar por que o gradiente é a \textbf{direção de máximo crescimento} (steepest ascent) e $-\nabla f$ é a direção de \textbf{máxima descida}.
    \item Discutir dificuldades práticas na \textbf{avaliação do gradiente} e aplicar \textbf{diferenças finitas} para aproximá-lo.
    \item Calcular a \textbf{taxa de variação de $f$ ao longo de uma direção} $S_i$:
    \[
      \frac{df}{d\lambda} = \nabla f^\top S_i.
    \]
    \item Descrever e aplicar o \textbf{método do declive máximo} (steepest descent / método de Cauchy) em um exemplo simples.
    \item Apresentar e interpretar \textbf{critérios de parada} em métodos iterativos de otimização irrestrita.
\end{itemize}

Baseado em \textit{RAO, Singiresu S. Engineering Optimization: Theory and Practice. 5. ed. Hoboken: John Wiley and Sons, 2020.}

\section{Contextualização}

\begin{itemize}
    \item Problema geral de \textbf{minimização irrestrita}:
    \[
      \min_{X \in \mathbb{R}^n} f(X).
    \]
    \item Nas aulas anteriores foram estudados métodos de \textbf{busca direta} (sem derivadas).
    \item Motivação para \textbf{métodos de busca indireta (descent methods)}:
    \begin{itemize}
        \item Usar informação de derivada para escolher direções de busca mais eficientes.
        \item Em problemas suaves, espera-se convergência mais rápida comparada a métodos puramente heurísticos.
    \end{itemize}
\end{itemize}

\section{Gradiente de uma função}

\subsection{Definição}

Seja $f:\mathbb{R}^n \to \mathbb{R}$ uma função escalar diferenciável. O \textbf{gradiente} de $f$ em $X = (x_1,\dots,x_n)^\top$ é o vetor coluna
\[
  \nabla f(X) =
  \begin{bmatrix}
    \dfrac{\partial f}{\partial x_1} \\[0.5em]
    \dfrac{\partial f}{\partial x_2} \\
    \vdots \\
    \dfrac{\partial f}{\partial x_n}
  \end{bmatrix}.
\]

\begin{itemize}
    \item $\nabla f(X)$ aponta, localmente, na direção em que $f$ cresce mais rapidamente.
    \item Geometricamente, em $\mathbb{R}^2$, o gradiente é ortogonal às curvas de nível de $f$.
\end{itemize}

\subsection{Intuição geométrica: direção de máxima inclinação}

\begin{itemize}
    \item Em $n=2$, imagine $f(x_1,x_2)$ como uma superfície.
    \item As curvas de nível $f(x_1,x_2)=c$ são linhas no plano $(x_1,x_2)$.
    \item O gradiente $\nabla f$ é perpendicular a essas curvas de nível e aponta para valores maiores de $f$.
    \item Consequentemente, $-\nabla f$ aponta para a direção de maior \textbf{decréscimo} de $f$.
\end{itemize}

\subsection{Teorema 6.3: o gradiente como direção de steepest ascent}

\textbf{Teorema 6.3.} \emph{O vetor gradiente $\nabla f(X)$ representa a direção de \textbf{máxima taxa de crescimento} (steepest ascent) da função $f$ no ponto $X$.}

\subsubsection*{Ideia da demonstração}

Considere um ponto arbitrário $X$ e um ponto vizinho $X + dX$ com
\[
  dX =
  \begin{bmatrix}
    dx_1 \\ dx_2 \\ \vdots \\ dx_n
  \end{bmatrix}.
\]

\begin{enumerate}[label=(\alph*)]
    \item O comprimento do vetor $dX$ é
    \[
      (ds)^2 = dX^\top dX = \sum_{i=1}^n (dx_i)^2.
    \]
    \item A variação de $f$ é dada por
    \[
      df = \sum_{i=1}^n \frac{\partial f}{\partial x_i} dx_i = \nabla f^\top dX.
    \]
    \item Escrevemos $dX = u\,ds$, onde $u$ é um vetor unitário ($\|u\|=1$) que indica a direção do passo.
    \item Então,
    \[
      \frac{df}{ds} = \nabla f^\top u.
    \]
    \item Usando o produto interno, obtemos
    \[
      \frac{df}{ds} = \|\nabla f\| \,\|u\| \cos\theta = \|\nabla f\| \cos\theta,
    \]
    onde $\theta$ é o ângulo entre $\nabla f$ e $u$.
    \item Como $\cos\theta \leq 1$, a taxa de variação $\dfrac{df}{ds}$ é máxima quando $\theta = 0^\circ$, isto é, quando $u$ está alinhado com $\nabla f$.
\end{enumerate}

\subsection{Teorema 6.4: valor máximo da taxa de mudança}

\textbf{Teorema 6.4.} \emph{A taxa máxima de variação de $f$ em um ponto $X$ é igual à norma do gradiente nesse ponto:}
\[
  \left(\frac{df}{ds}\right)_{\max} = \|\nabla f(X)\|.
\]

\noindent Isso decorre diretamente da expressão
\[
  \frac{df}{ds} = \|\nabla f\| \cos\theta,
\]
que atinge seu máximo em $\cos\theta = 1$.

\medskip

\noindent \textbf{Observação:} se o gradiente é a direção de máximo crescimento, então $-\nabla f$ é a direção de \textbf{máximo decréscimo} (steepest descent).

\section{Avaliação do gradiente na prática}

\subsection{Situações típicas problemáticas}

A avaliação de $\nabla f$ pode ser problemática em três situações:

\begin{enumerate}
    \item $f$ é diferenciável, mas é impraticável ou impossível derivar as expressões analíticas de $\dfrac{\partial f}{\partial x_i}$.
    \item As expressões analíticas existem, mas são muito complexas e caras de avaliar (alto custo computacional).
    \item O gradiente não é definido em todos os pontos (função não diferenciável em certos pontos).
\end{enumerate}

\noindent Nas situações (1) e (2), frequentemente se usam \textbf{diferenças finitas}. Na situação (3), o uso de diferenças finitas pode levar à quebra do algoritmo; métodos de busca direta são mais adequados.

\subsection{Diferenças finitas progressivas (forward difference)}

Para aproximar a derivada parcial em $X_m$:
\[
  \left.\frac{\partial f}{\partial x_i}\right|_{X_m}
  \approx
  \frac{ f(X_m + \Delta x_i u_i) - f(X_m) }{\Delta x_i},
\]
onde:
\begin{itemize}
    \item $\Delta x_i$ é um pequeno incremento na coordenada $x_i$;
    \item $u_i$ é o vetor canônico com $1$ na posição $i$ e $0$ nas demais.
\end{itemize}

\noindent Se $f(X_m)$ já for conhecido, essa aproximação requer apenas \textbf{uma avaliação adicional} de $f$ por componente. O gradiente completo em $\mathbb{R}^n$ requer $n$ avaliações adicionais de $f$.

\subsection{Diferenças finitas centrais (central difference)}

Uma aproximação mais precisa é dada por:
\[
  \left.\frac{\partial f}{\partial x_i}\right|_{X_m}
  \approx
  \frac{ f(X_m + \Delta x_i u_i) - f(X_m - \Delta x_i u_i) }{2\Delta x_i}.
\]

\noindent Nesse caso, são necessárias \textbf{duas avaliações} de $f$ por componente, em troca de maior precisão.

\subsection{Escolha de $\Delta x_i$: erros de arredondamento e truncamento}

\begin{itemize}
    \item Se $\Delta x_i$ é \textbf{muito pequeno}, os valores de $f$ em $X_m \pm \Delta x_i u_i$ são muito próximos e o erro de arredondamento pode dominar.
    \item Se $\Delta x_i$ é \textbf{muito grande}, perde-se a localidade da aproximação e o erro de truncamento aumenta.
    \item Na prática, escolhe-se $\Delta x_i$ de acordo com a escala de $x_i$ e a precisão numérica disponível.
\end{itemize}

\subsection{Gradiente não definido em todos os pontos}

Quando $f$ possui descontinuidades na derivada (quinas), a derivada pode não existir em certos pontos. Nesses casos:

\begin{itemize}
    \item Diferentes valores de $\Delta x_i$ podem produzir inclinações diferentes ao aplicar fórmulas de diferença finita.
    \item Isso é um indicativo de que a derivada não existe naquele ponto.
    \item Métodos baseados em gradiente podem falhar; é preferível usar métodos de busca direta.
\end{itemize}

\section{Taxa de variação ao longo de uma direção}

Em muitos métodos de otimização, escolhemos uma direção de busca $S_i$ e realizamos uma \textbf{busca unidirecional} em $\lambda$.

\subsection{Parametrização ao longo de $S_i$}

Dado um ponto $X_i$ e uma direção $S_i$, qualquer ponto sobre a reta é da forma
\[
  X(\lambda) = X_i + \lambda S_i.
\]
Definimos uma função escalar de uma variável:
\[
  \varphi(\lambda) = f(X_i + \lambda S_i).
\]

\subsection{Derivada direcional: $\dfrac{df}{d\lambda} = \nabla f^\top S_i$}

Pela regra da cadeia:
\[
  \frac{df}{d\lambda}
  = \sum_{j=1}^n \frac{\partial f}{\partial x_j} \frac{\partial x_j}{\partial \lambda}.
\]

Como $x_j(\lambda) = x_{ij} + \lambda s_{ij}$, temos $\dfrac{\partial x_j}{\partial \lambda} = s_{ij}$. Logo,
\[
  \frac{df}{d\lambda}
  = \sum_{j=1}^n \frac{\partial f}{\partial x_j} s_{ij}
  = \nabla f^\top S_i.
\]

\noindent Se $\lambda^\ast$ minimiza $\varphi(\lambda)$, então
\[
  \left.\frac{df}{d\lambda}\right|_{\lambda=\lambda^\ast} = 0
  \quad \Rightarrow \quad
  \nabla f(X_i + \lambda^\ast S_i)^\top S_i = 0.
\]

\noindent Esse resultado é fundamental em rotinas de \textbf{busca linear} (line search).

\section{Método do declive máximo (steepest descent / Cauchy)}

\subsection{Ideia básica}

\begin{itemize}
    \item Em cada iteração $i$, escolhemos a direção
    \[
      S_i = -\nabla f(X_i),
    \]
    que é a direção de \textbf{máxima descida} de $f$ a partir de $X_i$.
    \item Em seguida, realizamos uma busca unidimensional em $\lambda$ para encontrar o melhor passo $\lambda_i^\ast$.
    \item Atualizamos o ponto usando $X_{i+1} = X_i + \lambda_i^\ast S_i$.
\end{itemize}

\subsection{Algoritmo}

\begin{enumerate}
    \item \textbf{Inicialização:} escolher um ponto inicial $X_1 \in \mathbb{R}^n$ e definir $i = 1$.
    \item \textbf{Direção de busca:}
    \[
      S_i = -\nabla f(X_i).
    \]
    \item \textbf{Busca unidimensional:} determinar o passo ótimo
    \[
      \lambda_i^\ast = \arg\min_{\lambda \ge 0} f(X_i + \lambda S_i).
    \]
    \item \textbf{Atualização do ponto:}
    \[
      X_{i+1} = X_i + \lambda_i^\ast S_i = X_i - \lambda_i^\ast \nabla f(X_i).
    \]
    \item \textbf{Teste de convergência:} se algum critério de parada for satisfeito, encerrar; caso contrário, definir $i \leftarrow i+1$ e voltar ao passo 2.
\end{enumerate}

\noindent \textbf{Observação didática:} apesar de usar a melhor direção local de descida, o método pode ser \textbf{lento} em vales estreitos (comportamento de ``zigue-zague''), motivando o estudo de métodos mais avançados (Newton, quasi-Newton, gradiente conjugado, etc.).

\section{Exemplo: aplicação do método do declive máximo}

Considere
\[
  f(x_1,x_2) = x_1 - x_2 + 2x_1^2 + 2x_1x_2 + x_2^2.
\]

\subsection{Gradiente}

\[
  \nabla f =
  \begin{bmatrix}
    \dfrac{\partial f}{\partial x_1} \\[0.5em]
    \dfrac{\partial f}{\partial x_2}
  \end{bmatrix}
  =
  \begin{bmatrix}
    1 + 4x_1 + 2x_2 \\
    -1 + 2x_1 + 2x_2
  \end{bmatrix}.
\]

\subsection{Iteração 1}

\begin{itemize}
    \item Ponto inicial:
    \[
      X_1 =
      \begin{bmatrix}
        0 \\ 0
      \end{bmatrix}.
    \]
    \item Gradiente em $X_1$:
    \[
      \nabla f_1 = \nabla f(X_1) =
      \begin{bmatrix}
        1 \\ -1
      \end{bmatrix}.
    \]
    \item Direção de descida:
    \[
      S_1 = -\nabla f_1 =
      \begin{bmatrix}
        -1 \\ 1
      \end{bmatrix}.
    \]
    \item Ponto ao longo da direção:
    \[
      X(\lambda_1) = X_1 + \lambda_1 S_1 = (-\lambda_1, \lambda_1).
    \]
    \item Função ao longo da reta:
    \[
      f(-\lambda_1, \lambda_1) = \lambda_1^2 - 2\lambda_1.
    \]
    \item Derivada em relação a $\lambda_1$:
    \[
      \frac{df}{d\lambda_1} = 2\lambda_1 - 2.
    \]
    \item Passo ótimo:
    \[
      2\lambda_1 - 2 = 0 \;\Rightarrow\; \lambda_1^\ast = 1.
    \]
    \item Atualização:
    \[
      X_2 = X_1 + \lambda_1^\ast S_1 =
      \begin{bmatrix}
        -1 \\ 1
      \end{bmatrix}.
    \]
\end{itemize}

\subsection{Iterações seguintes (esboço)}

\begin{itemize}
    \item \textbf{Iteração 2:} calcular $\nabla f_2$, $S_2$, construir $f(X_2 + \lambda_2 S_2)$, encontrar $\lambda_2^\ast$, obter $X_3$.
    \item \textbf{Iteração 3 e seguintes:} repetir o procedimento até satisfazer um critério de parada.
    \item O ponto ótimo é
    \[
      X^\ast =
      \begin{bmatrix}
        -1.0 \\ 1.5
      \end{bmatrix}.
    \]
\end{itemize}

\noindent Em sala, pode-se:
\begin{itemize}
    \item Resolver a Iteração 2 com os alunos.
    \item Deixar as iterações seguintes como exercício.
    \item Discutir o comportamento de zigue-zague em curvas de nível.
\end{itemize}

\section{Critérios de parada}

Critérios usuais para encerrar o método iterativo:

\subsection{Variação relativa pequena em $f$}

\[
  \left|
  \frac{f(X_{i+1}) - f(X_i)}{f(X_i)}
  \right|
  \leq \varepsilon_1.
\]

\subsection{Norma do gradiente pequena}

\[
  \left| \frac{\partial f}{\partial x_j} \right| \leq \varepsilon_2,
  \quad j=1,\dots,n,
\]
ou, mais compactamente,
\[
  \|\nabla f(X_i)\| \leq \varepsilon_2.
\]

\subsection{Mudança pequena em $X$}

\[
  \|X_{i+1} - X_i\| \leq \varepsilon_3.
\]

\noindent Na prática, usualmente combina-se mais de um critério, com tolerâncias como $10^{-3}$, $10^{-4}$, etc., dependendo da aplicação.

\section{Encerramento e próximos passos}

\begin{itemize}
    \item Nesta aula:
    \begin{itemize}
        \item Definimos o gradiente e interpretamos seu papel como direção de \textbf{steepest ascent}.
        \item Vimos que $-\nabla f$ define uma direção de \textbf{steepest descent}.
        \item Apresentamos formas de \textbf{aproximar numericamente} o gradiente via diferenças finitas.
        \item Derivamos a taxa de variação de $f$ ao longo de uma direção $S_i$:
        \[
          \frac{df}{d\lambda} = \nabla f^\top S_i.
        \]
        \item Estudamos o \textbf{método do declive máximo} (Cauchy) e aplicamos a um exemplo.
        \item Discutimos \textbf{critérios de parada} em métodos iterativos.
    \end{itemize}
    \item Próximas aulas:
    \begin{itemize}
        \item Limitações do steepest descent em vales estreitos.
        \item Métodos mais avançados: \textbf{Newton}, \textbf{quasi-Newton}, \textbf{gradiente conjugado}, etc.
    \end{itemize}
\end{itemize}

\end{document}
